
\subsubsection{Model Compression and Optimization}

Reactive optimization techniques address large model constraints after discovering infeasibility. Quantization methods (Dettmers et al. 2022, 2023) reduce memory through lower precision representations (INT8, FP8, INT4), achieving up to 75\% memory reduction. Gradient checkpointing (Chen et al. 2016) trades computation for memory by recomputing activations during backward pass. Model parallelism frameworks including DeepSpeed (Rasley et al. 2020), Megatron-LM (Shoeybi et al. 2019), and FSDP (Zhao et al. 2023) distribute models across multiple GPUs.

These methods solve "how to make it fit" after discovering it does not fit. They require implementation effort to integrate and assume the configuration can be salvaged through optimization. Our work addresses a different question: validating feasibility before implementation begins. The two approaches are complementary—feasibility gates can identify configurations that either require reformulation or justify optimization investment before coding starts.

\subsubsection{Distributed Training and Infrastructure}

Production ML infrastructure addresses resource management through scheduling and allocation systems. Ray (Moritz et al. 2018) and Kubernetes-based systems enable dynamic resource allocation. vLLM (Kwon et al. 2023) optimizes inference serving. These systems assume infrastructure exists and focus on efficient utilization of runnable configurations. They do not validate whether a proposed experiment design is executable given available resources before implementation. Our proposed gate integrates into research workflows before infrastructure allocation decisions.

\subsubsection{Workflow and Pipeline Tools}

ML experiment management tools have matured significantly. MLflow (Zaharia et al. 2018), Weights \& Biases (Biewald 2020), and DVC handle versioning, reproducibility, and lineage tracking. Airflow and Kubeflow orchestrate complex pipelines. These tools focus on orchestration (workflow execution) and reproducibility (experiment tracking). Resource validation is limited to runtime monitoring—jobs fail when out-of-memory occurs—rather than proactive design-time checking. Our Phase 2C.5 gate adds an early checkpoint before orchestration begins.

\subsubsection{Negative Results Reporting}

The ML community increasingly recognizes negative results' value through dedicated tracks at NeurIPS and ICML, workshops on debugging ML models, and papers documenting hypothesis refutations (Lipton \& Steinhardt 2018), unexpected findings (Bender et al. 2021), and reproducibility challenges (Pineau et al. 2021). These typically report scientific findings (hypothesis refuted, method does not generalize) or implementation challenges (bugs, numerical instability). Computational infeasibility failures are rarely documented because they appear as project management failures rather than research contributions. We frame this as a workflow gap requiring process intervention.

\subsubsection{Positioning}

Our contribution is process-level. Where prior work provides tools to make large models run (compression, parallelism), manage running experiments (workflow tools), or document scientific failures (negative results), we identify a missing validation checkpoint in one automated research pipeline that resulted in wasted implementation effort. The feasibility gate sits between experiment design and implementation, proposing to complement reactive optimization with proactive validation.
